\subsubsection{Exact fermion construction used for the production data}
For all three Ising states, the production interval
matrices use the exact Jordan--Wigner solution. It is efficient
because a fermionic Gaussian state is specified by a $2N$ by $2N$
covariance matrix, rather than all $2^N$ amplitudes.
For the spin-to-fermion solution see Ref.~\cite{Montes2017}, Sec.~VIII
and Appendix C; for reduction by correlations see
Ref.~\cite{PeschelEisler}, Sec.~3.3(II), Eqs.~(14)--(18) and the
Majorana discussion immediately following them. The zero-mode parity
selection below fixes the conventions of our implementation explicitly.
\begin{enumerate}
\item \textbf{Transform spins to Majorana fermions.}
Define the Hermitian operators
\begin{equation}
 a_{2j}=\left(\prod_{k<j}Z_k\right)X_j,\qquad
 a_{2j+1}=\left(\prod_{k<j}Z_k\right)Y_j.
\end{equation}
They satisfy $a_ra_s+a_sa_r=2\delta_{rs}$.
In a fixed parity sector,
$H=(i/4)\sum_{r,s}\mathcal A_{rs}a_ra_s$, where the real
antisymmetric matrix has $\mathcal A_{r,r+1}=2$,
$\mathcal A_{r+1,r}=-2$, and boundary entry
$\mathcal A_{2N-1,0}=-2P$.
\item \textbf{Select the correct boundary sector.}
Even spin parity gives antiperiodic fermions, with momenta
$k=(2m+1)\pi/N$. Odd spin parity gives periodic fermions, with
$k=2m\pi/N$. In both cases $m=0,\ldots,N-1$, with momenta understood
modulo $2\pi$. These are two sectors of the same periodic \emph{spin}
Hamiltonian; we have not changed its physical boundary condition.
\item \textbf{Choose the occupation pattern for each primary.}
Diagonalizing the quadratic Hamiltonian defines fermion annihilation
operators $\gamma_k$ with excitation energies
$\omega_k=4|\sin(k/2)|$.
The identity state is their vacuum in the antiperiodic sector:
$\gamma_k\ket{I_N}=0$.
The energy primary occupies the two lowest modes in that sector:
\begin{equation}
 \ket{\eps_N}=\gamma_{\pi/N}^\dagger\gamma_{-\pi/N}^\dagger\ket{I_N}.
\end{equation}
For $\sigma$, use the periodic sector, empty its nonzero-energy modes,
and choose the zero-mode occupation that gives $P=-1$.
The zero mode must be fixed: a covariance with an unresolved zero mode
does not describe the required pure parity eigenstate.
\item \textbf{Construct the covariance in the implemented basis.}
Let $i\mathcal A=U\operatorname{diag}(\lambda_\mu)U^\dagger$.
For the identity, assign mode signs $\nu_\mu=\operatorname{sgn}\lambda_\mu$.
For $\eps$, reverse the signs for the two lowest positive-energy modes
and their negative-energy partners. Then set
\begin{equation}
 G=\operatorname{Re}\!\left[
 iU\operatorname{diag}(\nu_\mu)U^\dagger\right],
 \qquad G_{rs}=\langle ia_ra_s\rangle\quad(r\ne s),\quad G_{rr}=0.
\end{equation}
The mode signs are bookkeeping for the covariance, not occupation
probabilities. In the periodic sector start with $\nu_\mu=0$ on the
zero modes, take real orthonormal null vectors $u,v$ of $\mathcal A$,
and add either $uv^T-vu^T$ or its negative. Choose the sign such that
$(-1)^N\operatorname{Pf}(G)=-1$. The Pfaffian is the antisymmetric-matrix
polynomial obeying $\operatorname{Pf}(G)^2=\det G$; its sign here fixes
the physical fermion parity.
\item \textbf{Check purity and compute the interval matrix.}
Verify $G^T=-G$ and $G^2=-\id$. Compare the energy
$E=\frac14\sum_{r,s}\mathcal A_{rs}G_{rs}$ with
Eq.~\eqref{eq:ising-energies}.
Restrict the covariance to the interval and reconstruct its density
matrix using Eq.~\eqref{eq:gaussian}.
\end{enumerate}
This calculation is performed by \texttt{majorana\_covariance} and
\texttt{gaussian\_rdm}. The numerical purity and full-vector comparisons
are documented in Appendices~\ref{sec:state-comparisons} and
\ref{sec:large-sanity}.
